\section{Actual arbitrary-direction persistence gives source directional stability}

Write $v_e=1$ and $v_i=b_i$ for $i\ne e$, and $K$ for the actual whole
Taylor ideal of the original analytic point ideal $I$. All cone polynomials
below are the literal finite original homogeneous initial generators of $K$.

\begin{theorem}[Entire cone ray HS constancy from arbitrary-direction persistence]
\label{hir:actual-directional-persistent-cone-ray-hs}
\uses{hir:actual-analytic-point-normalized-bennett-all-orders,hir:actual-translated-initial-ray-family-isomorphisms,hir:actual-source-cone-polynomial-model}
\lean{Hironaka.HironakaActualPersistentDirectionalConeDerivatives.translated_cone_hs_eq_of_point_persistence}
\leanok
If the entire HS profile of the actual normalized point ideal in chart
$(e,b)$ equals that of $I$, then the actual translated cone ideal at $cv$
has that same entire profile for every complex $c$, including zero.
\end{theorem}
\begin{proof}
At $v$, the true normalized-to-translated-initial Bennett comparison and
the true translated-initial-to-original comparison sandwich every finite
HS value between equal values; hence equality holds. At every nonzero
multiple $cv$, the genuine homogeneous cone dilation is an actual original
germ algebra equivalence and preserves every finite quotient rank.
At zero, the whole original cone Taylor ideal is the whole homogeneous
initial ideal of $K$, whose entire formal HS profile equals that of $K$;
the actual original Taylor HS dictionary identifies this with the profile
of $I$. Thus the same equality holds along the ENTIRE ray.
\end{proof}

\begin{theorem}[Whole actual linear cone shear and tangent-axis derivative stability]
\label{hir:actual-directional-cone-linear-shear-axis-stability}
\uses{hir:actual-source-cone-polynomial-model,hir:actual-supplied-flattening-family-coupled-normal-data,hir:actual-centred-shear-whole-hs-invariance}
\lean{Hironaka.HironakaActualPersistentDirectionalConeDerivatives.polynomialShear,Hironaka.HironakaActualPersistentDirectionalConeDerivatives.polynomialShear_directional_chain,Hironaka.HironakaActualPersistentDirectionalConeDerivatives.homogeneous_polynomial_span_pderiv_of_line_hs}
\leanok
The actual invertible linear substitution $X_e\mapsto X_e$,
$X_i\mapsto X_i+b_iX_e$ retains homogeneous degrees and satisfies the
literal polynomial chain rule
$\partial_e(SP)=S(\delta_vP)$. For any finite homogeneous polynomial
family, constancy of the actual translated ideal HS along its pivot axis
implies $\partial_e$-stability of its ENTIRE formal span.
\end{theorem}
\begin{proof}
The linear substitution has inverse obtained by replacing $b$ with $-b$;
these identities hold pointwise on the actual complex vector space.
Actual composition therefore gives inverse original germ algebra maps.
On every coordinate generator, translating after substitution equals
substituting after translating by the actual linear image of the point.
Polynomial algebra-map extensionality proves this for the entire
polynomial, and mapping entire generator spans proves the whole translated
ideal equality. The actual germ equivalence preserves every finite HS
value. The chain rule follows on constants and variables from their literal
linear formulas, and on sums and products from the Leibniz rule.

To prove the tangent statement, swap the pivot with coordinate zero by an
actual coordinate linear equivalence. The translated HS equality becomes
constancy along the first coordinate plane. Genuine polynomial functions
are holomorphic on the whole actual ambient domain. The actual finite
Macaulay-rank dictionary converts the profile equality to every constant
rank along this plane. The actual supplied-family coupled-symbol theorem
then produces tangent derivative stability of the whole homogeneous
initial ideal of the translated origin ideal. Its actual Taylor ideal is
exactly the entire span of the original homogeneous polynomials; this span
is homogeneous, as seen by taking degree components in actual finite span
relations, so taking its homogeneous initial ideal changes nothing.
Undoing the actual pivot permutation gives the claimed entire-span
$\partial_e$-stability. No chart, rank, or derivative certificate is input.
\end{proof}

\begin{theorem}[Original arbitrary-direction persistence supplies the natural source derivation]
\label{hir:actual-original-directional-persistence-source-derivation}
\uses{hir:actual-directional-persistent-cone-ray-hs,hir:actual-directional-cone-linear-shear-axis-stability,hir:actual-reduced-canonical-homogeneous-initial-span}
\lean{Hironaka.HironakaActualPersistentDirectionalConeDerivatives.initial_directional_of_point_hs_persistence}
\leanok
For an arbitrary original analytic point ideal $I$ and actual point chart
$(e,b)$, entire actual source/current HS persistence implies
$\delta_vP\in\operatorname{in}K$ for every polynomial $P$ whose entire
formal image belongs to $\operatorname{in}K$.
\end{theorem}
\begin{proof}
Apply the first theorem to the actual original finite cone polynomials.
Apply the genuine linear shear to this SAME family, without choosing any
new standard basis. For a point $c\mathbf1_e$ on the pivot axis, the whole
translated-ideal shear identity identifies its HS profile with that of the
original cone at $cv$. The first theorem makes these all equal, so the
second theorem gives $\partial_e$-stability of the whole sheared span.
The actual polynomial-to-series substitution dictionary identifies that
span with $S(\operatorname{in}K)$, using genuine entire original cone
ideal generation. If $P\in\operatorname{in}K$, its genuine shear lies in
this entire sheared ideal. Its pivot derivative lies there also. The
literal polynomial chain rule identifies this derivative with
$S(\delta_vP)$, and the true inverse shear reflects whole ideal membership.
Thus $\delta_vP$ belongs to the original whole initial ideal. All
coefficients, source generators and original ideals are unchanged.
\end{proof}

This is a genuine point-chart source derivation producer. It does not itself
assert positive-dimensional-centre mixed maximal standardness or a history
termination theorem.

\begin{theorem}[Original arbitrary-direction persistence forces all same canonical point orders]
\label{hir:actual-original-arbitrary-point-persistent-canonical-orders}
\uses{hir:actual-original-directional-persistence-source-derivation,hir:actual-original-directional-persistent-canonical-orders-qualified}
\lean{Hironaka.HironakaActualPersistentDirectionalConeDerivatives.original_point_persistent_hs_orders}
\leanok
For any original analytic point ideal and actual point direction $(e,b)$,
actual entire source/current HS persistence forces every genuine quotient
of every SAME canonical reduced source series to retain its original vertex
order. No source stability qualification remains.
\end{theorem}
\begin{proof}
The previous theorem supplies exactly the natural source directional
polynomial derivation stability of the SAME whole source Taylor ideal.
The genuine directional canonical-order theorem then applies: actual
persistence internally supplies the whole current colons, and canonical
source sparsity plus the true coefficient bands exclude every low current
leading term. Neither this application nor the geometry chooses any new
source standard basis. It applies to ANY quotient satisfying the literal
normalization equation of the original canonical member.
\end{proof}

\begin{theorem}[One original canonical analytic family before the arbitrary-direction HS/order equivalence]
\label{hir:actual-original-arbitrary-point-hs-iff-orders-same-family}
\uses{hir:actual-original-arbitrary-point-persistent-canonical-orders,hir:actual-original-point-order-retention-entire-hs,hir:actual-original-canonical-reduced-vertex-basis,hir:actual-analytic-blowup-coordinate-pullback-taylor-naturality}
\lean{Hironaka.HironakaActualPersistentDirectionalConeDerivatives.exists_original_point_hs_iff_orders_family}
\leanok
For an arbitrary original analytic ideal $I$ in $\mathbb C^n$, a pivot
$e$ and arbitrary direction $b$, there exists ONE analytic family $H_a$
indexed by the SAME original canonical vertices, chosen before either
premise, such that
\[
 e^{|a|}H_a=\pi_{e,b}^*G_a,\qquad
 \mathcal N_{e,b}(I)=\langle H_a\rangle,
\]
where $\mathcal N_{e,b}(I)$ is the independently defined entire analytic
normalized-member ideal. For this SAME family,
\[
 \bigl(\forall\ell,\operatorname{HS}_{\mathcal N_{e,b}(I)}(\ell)
        =\operatorname{HS}_I(\ell)\bigr)
 \quad\Longleftrightarrow\quad
 \bigl(\forall a,H_a\in\mathfrak m^{|a|}\bigr).
\]
\end{theorem}
\begin{proof}
Choose the original canonical normalized analytic family by the accepted
original-only point generator theorem, before either side of the
equivalence. It gives the literal normalization equations, whole
independent ideal generation, and the forward order-retention-to-entire-HS
implication for this exact family. For the reverse implication, apply the
whole actual Taylor homomorphism to each normalization equation. Its
literal coordinate and analytic-blowup naturality gives the prescribed
formal quotient equation of the SAME canonical reduced source series.
The preceding original-only arbitrary-direction theorem puts this Taylor
series in its original vertex power. The genuine whole original Taylor
point-power contraction reflects this to the actual analytic maximal
power of $H_a$. Dimension reassociation between $0+n$ and $n$ transports
the prescribed pivot and each coefficient of the SAME direction explicitly;
it makes no new coordinate choice and no family substitution. This proves
both directions for the family already chosen before the hypotheses.
\end{proof}
